% Transcription — fonds Alexandre Grothendieck, shelfmark 156-7, batch 1 (pages 1-20).
% Established from the facsimile archives/batches/156-7/batch-01.pdf (a copy of
% 156-7.pdf fetched through the project's relay; no cover sheet in the batch,
% so PDF page k = folder page k), per the skill transcribe-grothendieck. The
% folder is 113 pages in six batches (1-20, 21-40, 41-60, 61-80, 81-100,
% 101-113), transcribed in parallel; this is the first. The notation follows
% the second draft, folder 156-6 (GF VI), whose headers were read first.
%
% Pass: Opus 5.5 (claude-opus-5-5), 2026-09-26 — first pass, unchecked against the pages by a human.
%
% Pages skipped: none. Pages summarised: none. All twenty leaves are
% mathematics in his hand and all are transcribed.
%
% His pagination 1-20 runs at the head of the leaves and coincides with the
% archivists' (noted once, on page 1; the « 4 » is overwritten). Page 1 carries
% the heading « Analysis situs (troisième mouture) », top right in a ruled
% corner « GF VII », and below it a slanted cross-reference « cf. ... p. 23,
% 50-58, 59, 110-113 » (the last numbers circled). The one date in his hand
% in this batch is « 23 juin », underlined, opening page 1 (set as
% \subsection{23 juin} with a \note). The bracketed « [Chapitre] » of the
% inventory title is the archivists'.
%
% Numbered runs. Preliminary data lettered by small circled signs, read
% (a)-(d) on page 3 and (e)-(f) on page 4 (the letters are doubtful). Axioms:
% At 0 (p. 3, unnumbered by his own account); « I) Axiomes sine qua non »
% At 1 - At 7 (pp. 4-5), with At 5 / At 6 renumbered by overwriting on page 5
% (the page gives At 4, At 6, then At 5, At 7) and At_M 4 - At_M 6 on page 6
% swapped by a curved arrow; At_M 0 - At_M 3 (p. 5), At_M 4 a), b), At_M 5,
% At_M 6 (p. 6); At DC 1 (p. 7, cancelled), At CD 1, At CD 1'_0, At CD 2,
% At CD 2' (pp. 8-10); At pol 1 twice (p. 13, the good-intersection condition;
% p. 16, a compatibility criterion — the number overwritten), At simp 1 (p. 14);
% « axiomes complémentaires » At C, At D (pp. 14-15), At D_0, At D' (pp. 16-17),
% At spéc, At fil (p. 17), At supp, At ens D (p. 18); axioms on lieux
% At L 1 - At L 3 and At CL (pp. 19-20). Statements: Propriété 1, 1'_0, 2
% (pp. 8-9), Cor. 1-3 (p. 9, and Corollaire 4 in the margin), Cor. Main 1, 2
% (p. 8), Cor. Main 1 again (pp. 16, 17), Corollaire 2 (p. 18), Cor 1-3
% (p. 19), two unnumbered Cor (p. 20). No formula numbers in this batch.
%
% Where this draft restates the second one (GF VI = 156-6), noted where it
% occurs: the explicit cross-references « GF p. 28 » and « p. 32 » (p. 6, the
% Scholie and Théorème-Scholie on At*1-At*5), « cf. GF VI, p. 73 » (p. 8,
% struck, with the Proposition that GF VI turned into At 12), « At 11 p. 67,
% 68 » (p. 8 margin), « GF VI ... p. 52 » (p. 12 margin), and « GF V p. 19,
% 20 » (p. 12, chapter V = 156-5). Also noted without his reference: the
% smallest-stratum axiom (GF VI At 5, here At 6), recollement des raffinements
% (GF VI At 11, here Cor. Main 1 of At C, p. 17), At L 1 = GF VI At 8 (p. 19),
% At L 2 = GF VI At 9 (p. 20).
%
% What the batch is about. The third draft opens (23 June) by recapitulating
% the purely « algebraic » data and axioms of an atelier: figures 𝔉 with ≤
% (sub-figure) and ≪ (refinement), At 0 (≤ ⇒ ≪), the operations F ∨ G,
% F ∧ G, F·G = Inf^≪, ∅_𝔉, multistrates 𝓜 (strictly irreducible figures),
% lieux ℒ (≪-minimal non-empty figures); déploiement F̃, grande/petite
% (multi)ombre. Pages 4-6: axioms At 1-At 7 (sums, strata, compatibility by
% strata, refinement by strata, smallest stratum, interior refinement ≪̊ and
% its transitivity), a Scholie reading them in (𝓜, ≤, ≪, 𝔉̃), and the
% reduction to (𝓜, ≤, compatibility) for figures of finite type. Pages 7-12:
% disjointness ∥ and the compatibility/disjunction axioms At CD 1, 1'_0, 2, 2'
% (ateliers « spécieux » and « polyédraux »), with a hesitation over ≪ for
% simplices (a Δ¹ inside a Δ², subdivisions of Δ² that nobody would call
% simplicial). Pages 13-14: polyhedral and simplicial ordered sets and
% ateliers (At pol 1, At simp 1). Pages 14-18: « axiomes complémentaires »
% At C, At D (compatibility and disjunction of refinements of sub-figures,
% with the Omb diagram of three cartesian squares on page 15), At pol 1 /
% At spéc, At fil (filtering Sups), At supp, At ens D. Pages 19-20: the
% axioms on lieux — At L 1 (divisibility: omb(X)° ≠ ∅, « l'infini actuel
% tout court ! »), At L 2, At L 3, At CL; page 20 breaks off mid-sentence.
%
% Notation, as in 156-6: 𝔉 \mathfrak{F}, 𝓜 \mathcal{M}, ℒ \mathcal{L}; his
% crossed chevrons (compatibility) \between; disjointness \parallel, barred
% \overline{\parallel}; incidence \triangleleft; interior refinement
% \overset{\circ}{\ll}; « Omb », « omb », « Multomb », « Fig », « Sup »,
% « Inf » in roman. His looped capital for a family of parts is \Phi (and
% \Psi), as in 156-6. The axiom labels are spelt as he writes them here: « At 1 »
% ..., « At_M » (with 𝓜 as index), « At CD », « At pol », « At simp » (second
% word doubtful), « At spéc », « At fil », « At supp », « At ens D »,
% « At L », « At CL ». Drawings: the Omb diagram (p. 15, headless lines) and
% the déploiement diagram (p. 15) and the small scheme of p. 17 are set in
% tikz-cd as drawn; the diamond-and-circle of p. 8 and the triangles of p. 11
% are described in French notes. His underlining is \emph.
%
% Hard pages: 2 (a long slanted margin note, read in fragments), 6 and 8
% (interlinear rewriting of At_M 4 and At CD 1, cancelled blocks), 12-13
% (cancelled passages and a long slanted margin), 16 (heavily struck lines
% and a margin partly cancelled), 18 (long slanted margin note), 20 (At CL,
% much retouched). Apparatus: 148 \ill{} and 86 \uncertain{}. Variances on
% the page kept and flagged: « Inf^≪(X, L) » for X' (p. 14); K introduced and
% L used in At CD 1, At CD 2 (pp. 8, 10); two different statements called
% At pol 1 (pp. 13, 16); the At 5 / At 6 order (p. 5).
%
\documentclass[11pt,a4paper]{article}
\input{../preamble/grothendieck}

\folder{156-7}
\batch{1}
\pages{1}{20}
\dating{1986}
\watermark{Édition de démonstration}
\foldertitle{[Chapitre] VII. Analysis situs (troisième mouture) : notes manuscrites (23-26/06/1986)}

\begin{document}
\page{1}
\section{Analysis situs (troisième mouture)}

\note{titre de sa main en tête de la page 1, qui porte en haut au centre son numéro « 1 » ; ses numéros de page, en tête de chaque feuillet, coïncident dans ce lot avec la numérotation des archivistes. Dans l'angle supérieur droit, dans un coin tracé à la plume, « GF VII » ; au-dessous, en oblique, un renvoi de sa main : « cf.\ \ill{} p.\ 23, 50-58, 59, 110-113 », les derniers nombres cerclés (le « 50 » est écrit en surcharge, et un mot au-dessus des nombres n'est pas lu)}

\subsection{23 juin}
\note{date de sa main, soulignée, en tête du premier alinéa}

Je récapitule les données et axiomes d'une « algèbre de figures », ou « \emph{atelier} ». Il s'agit pour l'instant d'axiomes purement « algébriques » \uncertain{dominant} le formalisme algébrique des figures, à l'exclusion pour le moment de propriétés topologiques (dimension, connexité, acyclicité) ou \emph{proprement} géométriques (régularité topologique, boules et sphères topologiques d'une algèbre de figures), voire les « propriétés combinatoires finies », \uncertain{tournant autour} de la notion de \uncertain{simplexe}.

Parmi les axiomes, je distingue les axiomes « \uncertain{sine qua non} » At 1 à At 7, qui me paraissent les plus essentiels pour fonder la \uncertain{seule} géométrie topologique du formalisme envisagé. Les autres axiomes seront considérés comme plus ou moins « facultatifs » au niveau du formalisme algébrique. Ils pourront avoir un rôle crucial, dès qu'il s'agira d'un travail sur pièces dans le

\page{2}
contexte d'un atelier définissant une « contrée ».

Le travail précédent a permis de dégager le point de vue le plus commode pour le formalisme algébrique. La notion fondamentale me semble bien celle de \emph{figure}, et non pas celle de « figure élémentaire » ou « multistrate », voire celle de « lieu » (voire, \ill{} la notion de « lieu ponctuel » — ou « point » \ill{} — qui \ill{} par avoir été \ill{}…). Une des raisons en est la \uncertain{simplicité} du formalisme algébrique, \uncertain{quand} il est agréable de pouvoir manipuler les figures et multistrates comme étant des figures particulières. Quand une figure $F$ est interprétée \add{(par l'ens.\ de ses « strates »)} comme une partie \add{(ordonnée par $\leq$)} de l'ens.\ $\mathcal{M}$ des multistrates, nous allons écrire $\widetilde{F}$ (et non $F$) pour cette partie. Quand elle est interprétée comme une figure \add{ensembliste} dans $\mathcal{M}$, nous l'écrivons $\mathrm{Multomb}(F)$ (« grande multiombre »), son support étant $\mathrm{Omb}(F)$ (la « grande ombre » de $F$) — cette dernière, contrairement à $\widetilde{F}$ et à $\mathrm{Multomb}(F)$, ne permettant pratiquement jamais de reconstituer $F$, mais moralement, son support « seulement

\marginal{L'autre raison \ill{} « \ill{} » \ill{} simplicité \ill{} figures \ill{} multistrates}
\note{longue note oblique dans la marge gauche, sur une quinzaine de lignes courtes, lue seulement par fragments}

\page{3}
» (et je suis tenté \uncertain{maintenant} sur ce support dans la mouture précédente GF VI). Quand on interprète $F$ comme une figure ensembliste dans l'ens.\ des lieux $\mathcal{L}$, \struck{ce} (pratiquement ceci signifie que c'est un atelier ensembliste) on l'écrira $\mathrm{multomb}(F)$ (petite multiombre), son support sera $\mathrm{omb}(F)$.
\note{la parenthèse ouverte au bas de la page 2 (« seulement ») se ferme au haut de la page 3 ; la lecture de la première ligne est incertaine}

\emph{Données} \emph{préliminaires} :
\begin{itemize}
  \item[$\mathfrak{F}$] ensemble de figures
  \item[$\leq$] relation d'\emph{inclusion} des figures, $F \leq G \overset{\mathrm{déf}}{\Longleftrightarrow} F$ sous-figure de $G$
  \item[$\ll$] relation de raffinement, $F \ll G \overset{\mathrm{déf}}{\Longleftrightarrow} F$ raffine $G$
\end{itemize}
\note{les trois lignes sont réunies par une accolade à gauche}

\emph{At 0} \quad avec
\[
  F \leq G \Longrightarrow F \ll G
\]
— axiome si évident que je ne \uncertain{le mets} pas \uncertain{inclus} dans les axiomes numérotés.

\marginal{NB.\ Ce \uncertain{serait} \ill{} dans At 4 plus bas}
\note{note de sa main dans la marge gauche, en face de « At 0 »}

\begin{itemize}
  \item[(a)] $F \vee G$ ou $F \cup G$ \quad figure « engendrée » ou « réunion » ou « somme » : c'est le sup au sens de $\leq$ (quand il existe) — auquel cas c'est aussi un sup pour $\ll$ \add{ou Inf pour $\leq$} ;
  \item[(b)] $F \wedge G$ ou $F \cap G$ ($= \mathrm{Inf}^{\leq}(F, G)$) \quad figure \emph{intersection} (existe toujours \uncertain{(même pour familles infinies (non vides))} ; quand \uncertain{en outre} $F \vee G$ existe, alors c'est aussi le Inf pour $\ll$ ;
  \item[(c)] $F \cdot G$ ($= \mathrm{Inf}^{\ll}(F, G)$) \quad figure « intersection \emph{fine} » (quand elle existe) : Inf pour $\ll$ ;
  \item[(d)] $\varnothing_{\mathfrak{F}}$ \quad figure vide : plus petit élément de $\mathfrak{F}$ pour $\leq$ \emph{et} pour $\ll$.
\end{itemize}
\note{chacune des quatre lignes est précédée d'un petit signe cerclé, lu ici (a) à (d) sans certitude ; sous « $F \cdot G$ » un mot court est biffé. L'ajout « ou Inf pour $\leq$ » de (a) est écrit plus bas dans la colonne de droite et rattaché par un trait ; sa place n'est pas sûre}

\marginal{NB : Si $(F, G)$ ont un sup pour $\ll$ (c'est toujours le cas dans le cas ensembliste $\mathrm{Fig}(\mathcal{L})$), ce n'est pas nécessairement un sup pour $\leq$}
\marginal{NB \quad $\widetilde{F \cap G} = \widetilde{F} \cap \widetilde{G}$}
\marginal{NB Existe dans le cas ensembliste $\mathrm{Fig}(\mathcal{L})$ et dans le cas \uncertain{rectiligne}}
\note{trois notes de sa main dans la colonne de gauche, en face de (a)-(b) et de (c) ; la première est ouverte par un mot biffé (peut-être « Question »)}

\page{4}
\begin{itemize}
  \item[(e)] $\mathcal{M} =$ ensemble des « \emph{figures} \emph{élémentaires} » ou \emph{multistrates} $\subset \mathfrak{F}$ \quad formé des objets de $(\mathfrak{F}, \leq)$ strictement irréductibles. Muni des relations $\underset{\mathcal{M}}{\leq}$, $\underset{\mathcal{M}}{\ll}$ induites par $\leq$ et $\ll$ dans $\mathfrak{F}$ ;
  \item[(f)] $\mathcal{L} \subset \mathcal{M}$, ens.\ des « \emph{lieux} » : éléments minimaux de $\mathfrak{F} \smallsetminus \lbrace \varnothing_{\mathfrak{F}} \rbrace$ pour $\ll$.
\end{itemize}
\note{les deux lignes sont précédées, comme celles de la page 3, d'un petit signe cerclé ; au-dessus de la première, un chiffre surchargé. « lieux » est écrit au-dessus d'un mot souligné et biffé ; la définition des lieux, dans la colonne de droite, est rattachée à (f) par un long trait}

\marginal{strates d'une figure : les $X \in \mathcal{M}$ tels que $X \leq F$ ; \emph{déploiement} de $F$ : ens.\ de ses strates $\widetilde{F} = \lbrace X \in \mathcal{M} \mid X \leq F \rbrace$. On écrit $X \triangleleft F$ ($X$ \emph{incident} à $F$) pour $X \in \mathcal{M}$ et $X \leq F$}
\note{note de sa main dans la marge gauche, encadrée d'un trait qui la relie à (e)}

\emph{Axiomes}

I) Axiomes « sine qua non ». Les premiers trois \add{concernent} (At 1 – At 3) \struck{\ill{}} $(\mathfrak{F}, \leq)$, les quatre suivants At 4 – At 7 les relations entre $\leq$ et $\ll$.
\note{« At 1 – At 3 » est écrit à droite, au-dessus de la ligne ; un « At 4 » est biffé au début de la ligne suivante}

\emph{At 1} \quad $\mathfrak{F} \neq \varnothing$ et stable par $\mathrm{Sup}_{I}$ majorés ($I$ quelconque)

\marginal{Axiome des sommes majorées de figures}

\emph{At 2} \quad Toute figure est le sup de ses strates :
\[
  F = \mathop{\mathrm{Sup}}_{X \in \widetilde{F}} X \qquad \text{où} \qquad \widetilde{F} = \lbrace X \in \mathcal{M} \mid X \leq F \rbrace
\]

\marginal{Axiome des strates}
\marginal{équivalent à : si $F, G \in \mathfrak{F}$, alors $F \leq G \Leftrightarrow \forall\, X \in \widetilde{F}$, $\exists\, Y \in \widetilde{G}$ avec $X \leq Y$ — \emph{Critère d'inclusion des figures par strates}}
\note{notes de sa main en oblique dans la marge gauche, en face de At 1 et At 2}

\emph{Scholie} \quad At 1 et At 2 signifient que $\mathfrak{F}$ s'interprète, via $F \mapsto \widetilde{F}$, comme une partie $\widetilde{\mathfrak{F}}$ de $\mathfrak{F}(\mathcal{M})$, $\mathcal{M}$ ens.\ ordonné (par $\leq$), avec les conditions :
\begin{itemize}
  \item[a)] les $\Phi \in \widetilde{\mathfrak{F}}$ sont \uncertain{fermées}
  \item[b)] si $\Phi \in \widetilde{\mathfrak{F}}$, toute partie fermée de $\Phi$ est $\in \widetilde{\mathfrak{F}}$
  \item[c)] \struck{$\widetilde{\mathfrak{F}} \neq \varnothing$ ($\Leftrightarrow \varnothing \in \widetilde{\mathfrak{F}}$)} $\bigcup_{\Phi \in \widetilde{\mathfrak{F}}} \Phi = \mathcal{M}$, et $\widetilde{\mathfrak{F}} \neq \varnothing$
\end{itemize}
\note{« $\mathfrak{F}(\mathcal{M})$ » ainsi lu ; sous le c) biffé, un mot entre parenthèses n'est pas lu}

ou, indépendamment d'une relation d'ordre sur $\mathcal{M}$, un ens.\ de parties satisfaisant
\begin{itemize}
  \item[b')] $\forall\, X \in \mathcal{M}$, l'ens.\ $\widetilde{\mathfrak{F}}^{X} = \lbrace \Phi \in \widetilde{\mathfrak{F}} \mid X \in \Phi \rbrace$ a un plus petit élément (noté $\widetilde{X}$)
  \item[d')] Dans $(\widetilde{\mathfrak{F}}, \subset)$, toute famille majorée a un Sup, \emph{et} $\widetilde{\mathfrak{F}} \neq \varnothing$ ($\Leftrightarrow \varnothing \in \widetilde{\mathfrak{F}}$)
\end{itemize}
\note{b') et d') sont réunis à gauche par une accolade ; au début de d') un « $\forall$ » est biffé}

\marginal{$\widetilde{\mathfrak{F}} \neq \varnothing$ \ill{} \ill{} et donc $\varnothing \in \widetilde{\mathfrak{F}}$}
\note{note oblique dans la marge gauche, en face de c), lue en partie ; une flèche la mène vers c)}

\emph{Relation}
\[
  F \between G \overset{\mathrm{déf}}{\Longleftrightarrow} F \vee G \text{ existe, i.e.\ } \lbrace F, G \rbrace \text{ majoré dans } (\mathfrak{F}, \leq)
\]
\emph{figures compatibles}. Cette relation induit $X \between Y$ dans $\mathcal{M}$. Dans le cas où les figures sont de t.f., alors $(\mathfrak{F}, \leq)$ \add{avec (At 1, At 2, At 3)} \struck{\ill{}} se reconstitue par $(\mathcal{M}, \leq, \between)$, où $\between$ satisfait
\[
  \begin{cases}
    X \between X & \text{réflexive, symétrique} \\
    X \between Y,\ X' \leq X,\ Y' \leq Y \Longrightarrow X' \between Y'
  \end{cases}
\]
\note{« t.f. » : de type fini}

\page{5}
Donc la donnée de $(\mathfrak{F}, \leq)$ avec figures de t.f.\ équivaut à : ($\mathcal{M}$, $\leq$ \struck{\ill{}} (ordre), $\between$ (relation réflexive symétrique)) satisfaisant
\begin{itemize}
  \item[a)] \struck{\ill{}} \add{pour} Toute partie \add{fermée} majorée de $\mathcal{M}$ est de type fini (i.e.\ réunion finie d'ens.\ $\widetilde{X} = \mathcal{M}_{\leq X}$)
  \item[b)] $X \between Y$, $X' \leq X$, $Y' \leq Y \Longrightarrow X' \between Y'$ !
\end{itemize}
\note{au début de a), une formule de deux lignes est lourdement biffée}

\marginal{On aura alors $\widetilde{\mathfrak{F}} =$ ens.\ des parties fermées de $\mathcal{M}$, de type fini, formées d'éléments deux à deux compatibles (pour $\between$)}

\emph{At 3} \quad Soient $F, G \in \mathfrak{F}$ tels que $\forall\, X \in \widetilde{F}$, $Y \in \widetilde{G}$, on ait $X \between Y$, alors $F \between G$ (i.e.\ $F \vee G$ existe)

\marginal{Critère de compatibilité par strates}

\struck{At} Les axiomes suivants concernent les relations entre $\leq$ et $\ll$.

\emph{At 4} \quad $\forall\, F, G \in \mathfrak{F}$,
\[
  F \ll G \Longleftrightarrow \forall\, X \in \widetilde{F},\ \exists\, Y \in \widetilde{G},\ \text{t.q.\ } X \ll Y
\]
\uncertain{Ramène} la connaissance de $\ll$ dans $\mathfrak{F}$ à celle de $\underset{\mathcal{M}}{\ll}$.

\marginal{Critère de raffinement par strates [Corollaire $F \leq G \Longrightarrow F \ll G$]}

\emph{At 6} \quad \struck{$\forall$} $X, F \in \mathfrak{F}$, $X \triangleleft F$ (i.e.\ $X \in \widetilde{F}$), soit $F^{X} = \lbrace Y \in \widetilde{F} \mid X \ll Y \rbrace$. Alors $F^{X}$ a un plus petit élément.
\note{le chiffre de « At 6 » est surchargé (un 6 sur un 5, semble-t-il) ; dans la deuxième mouture, l'axiome de la plus petite strate est At 5 (GF VI, p.\ 21-22)}

\marginal{Axiome de la plus petite strate, ou du raffinement interne}

\emph{Notation} \quad $X \mathrel{\overset{\circ}{\ll}} Y \overset{\mathrm{déf}}{\Longleftrightarrow} X$ est plus petit élément de $\widetilde{Y}^{X}$, i.e.\ $X \ll Y$, et $X \ll Y' \triangleleft Y \Rightarrow Y' = Y$ (ou $X \in \mathrm{Omb}(Y)^{\circ}$) — $X$ \emph{raffinement interne} de $Y$.

\emph{At 5} \quad Si $X \between Y$ ($X, Y \in \mathcal{M}$) alors $X \ll Y \Rightarrow X \leq Y$
\note{le 5 est écrit en surcharge, très appuyé, sur un autre chiffre (6 ?) ; au-dessus de « Si », le mot « figures » est ajouté et relié par un trait à la note marginale}

on en conclut, dans une $\widetilde{F}$, $\leq$ et $\ll$ coïncident, \uncertain{ou encore}, si $F, G \in \mathcal{M}$ et $F \between G$, alors $F \ll G \Leftrightarrow F \leq G$.

\marginal{équivalent à : si $G, H$ \ill{} des figures \ill{} \ill{} $F \ll G \Rightarrow F \leq G$ — Critère des raffinements \uncertain{autonomes}}

\emph{At 7} \quad Soient $X, Y, Z \in \mathcal{M}$, alors $X \mathrel{\overset{\circ}{\ll}} Y \mathrel{\overset{\circ}{\ll}} Z \Longrightarrow X \mathrel{\overset{\circ}{\ll}} Z$
\note{le 7 est écrit en surcharge sur un autre chiffre}

\marginal{Transitivité des raffinements intérieurs}
\note{les notes marginales de cette page, écrites en oblique dans la marge gauche en face de chaque axiome, sont soulignées et reliées aux axiomes par des traits}

Traduction en termes de $(\mathcal{M}, \leq, \ll, \widetilde{\mathfrak{F}})$
\begin{itemize}
  \item[$\mathrm{At}_{\mathcal{M}}0$] $\leq \Longrightarrow \ll$
  \item[$\mathrm{At}_{\mathcal{M}}1$] Les $\Phi \in \widetilde{\mathfrak{F}}$ fermés pour $\leq$
  \item[$\mathrm{At}_{\mathcal{M}}2$] Si $\Phi \in \widetilde{\mathfrak{F}}$, \struck{toute partie} \add{$\Psi$} fermée de $\Phi$ pour $\leq$, alors $\Psi \in \widetilde{\mathfrak{F}}$, et $\widetilde{\mathfrak{F}} \neq \varnothing$ (i.e.\ $\varnothing \in \widetilde{\mathfrak{F}}$)
  \item[$\mathrm{At}_{\mathcal{M}}3$] Si $\Phi, \Psi \in \widetilde{\mathfrak{F}}$ telles que $\forall\, X \in \Phi$, $\forall\, Y \in \Psi$, \struck{\ill{}} $X \between Y$ (i.e.\ $\exists\, \Phi \in \widetilde{\mathfrak{F}}$, $X, Y \in \Phi$), alors $\Phi \between \Psi$ (i.e.\ $\Phi \cup \Psi \in \widetilde{\mathfrak{F}}$).
\end{itemize}
\note{« $\Phi$ » rend sa capitale bouclée, comme dans la deuxième mouture ; dans la dernière parenthèse, le signe de compatibilité porte un $\mathcal{M}$ en indice}

\page{6}
\struck{$\mathrm{At}_{\mathcal{M}}4$ : $\forall\, X \in \mathcal{M}$, $\Phi \in \widetilde{\mathfrak{F}}$, posant $\Phi^{X} = \lbrace Y \in \Phi \mid X \ll Y \rbrace$, $\Phi^{X}$ a un plus petit élément}

$\mathrm{At}_{\mathcal{M}}4$ \quad a) \struck{Posons $\forall\, X, Y \in \mathcal{M}$, $X \mathrel{\overset{\circ}{\ll}} Y \overset{\mathrm{déf}}{\Longleftrightarrow} X$ \ill{} minimal} Si $X \ll Y$, alors il y a un \struck{\ill{}} \add{minimal} dans $\widetilde{Y}^{X} = \lbrace Z \in \widetilde{Y} \mid X \ll Z \rbrace$ \struck{posons $X \mathrel{\overset{\circ}{\ll}} Y$ si \ill{} i.e.\ $X \ll Y' \leq Y \Rightarrow Y' = Y$}

b) Alors $X \mathrel{\overset{\circ}{\ll}} Y$, $X \mathrel{\overset{\circ}{\ll}} Y'$, $Y \between Y' \Longrightarrow Y = Y'$
\note{a) est un tissu de ratures et d'ajouts interlinéaires, dont l'ordre n'est pas sûr}

$\mathrm{At}_{\mathcal{M}}6$ \quad Si $X \between Y$, alors $X \ll Y \Longrightarrow X \leq Y$

$\mathrm{At}_{\mathcal{M}}5$ \quad $X \mathrel{\overset{\circ}{\ll}} Y \mathrel{\overset{\circ}{\ll}} Z \Longrightarrow X \mathrel{\overset{\circ}{\ll}} Z$
\note{les chiffres 6 et 5 sont écrits en surcharge et les deux lignes sont reliées par une flèche courbe qui les intervertit ; la numérotation At 5, At 6 de la page 5 et celle-ci ne se correspondent pas terme à terme, et elles restent telles quelles}

\marginal{J'ai exprimé At 5 \ill{} \ill{} $\mathfrak{F}$ \ill{} $\widetilde{\mathfrak{F}}$, mais directement en termes de $(\mathcal{M}, \leq, \ll, \between)$}
\note{note de sa main dans la marge gauche, avec une flèche vers « $\mathrm{At}_{\mathcal{M}}$ » en haut de la page et une autre vers $\mathrm{At}_{\mathcal{M}}6$}

D'autre part, pour des $\mathfrak{F}$ à figures de type fini, \add{bien \uncertain{connues}} \ill{} l'élément de structure $\widetilde{\mathfrak{F}}$ \add{\ill{} au profit de $\between$}, et remplacer $\mathrm{At}_{\mathcal{M}}1$ à $\mathrm{At}_{\mathcal{M}}3$ par (a) et la

\struck{$\mathrm{At}_{\mathcal{M}}$(\ill{}) Toute partie fermée majorée de $\mathcal{M}$ est de t.f.}

\struck{\ill{}} \uncertain{Voir} p.\ 5 (Toute partie fermée majorée de $\mathcal{M}$ est de t.f.) et la condition b sur $\between$. (page 5)

\marginal{i.e.\ $X \between Y$ et $X' \leq X$, $Y' \leq Y \Longrightarrow X' \between Y'$}
\note{note oblique dans la marge gauche, en face de ces lignes ; elle redit la condition b) de la page 5}

Dans le cas général, grâce à \uncertain{laquelle} $\between$ satisfait cette condition b) plus les conditions $\mathrm{At}_{\mathcal{M}}4, 5, 6$, on trouve
\[
  \widetilde{\mathfrak{F}}_{\mathrm{tf}} = \text{ensemble des parties fermées } \Phi \text{ de t.f.\ de } (\mathcal{M}, \leq) \text{ telles que } X, Y \in \Phi \Rightarrow X \between Y .
\]
\note{« $\Phi$ de t.f. » est ajouté au-dessus de la ligne ; « telles que » est écrit au-dessus d'un mot biffé}

Pour avoir $\widetilde{\mathfrak{F}}$, il faut cependant un choix de structure en plus : se donner certains antifiltres infinis dans l'ens.\ ordonné $\widetilde{\mathfrak{F}}_{\mathrm{tf}}$ précédent…

On peut aussi interpréter $\mathfrak{F}$ en termes de
\[
  \mathfrak{F}^{*} = \lbrace \mathrm{Multomb}(F) \mid F \in \mathfrak{F} \rbrace \subset \mathrm{Fig}(\mathcal{M})
\]
où
\[
  \mathrm{Multomb}(F) = \lbrace \mathrm{Omb}(X) \mid X \in \widetilde{F} \rbrace , \qquad
  \mathrm{Omb}\, X = \lbrace Y \in \mathcal{M} \mid Y \ll X \rbrace = \mathcal{M}_{\ll X} .
\]
(cf.\ Scholie GF p.\ 28, et axiomes, en \uncertain{nombre} de \ill{} p.\ 32), comme un sous-atelier (pas néc.\ spécieux) de $\mathrm{Fig}(\mathcal{M})$.
\note{dans $\mathrm{Multomb}(F)$, un signe est biffé devant « Omb ». Les renvois sont à la deuxième mouture : GF VI, p.\ 28 porte le Scholie décrivant un atelier par $\mathcal{M}$ et une famille $\mathfrak{F}_{\mathcal{M}} \subset \mathrm{Fig}(\mathcal{M})$, et p.\ 32 le Théorème-Scholie sur les conditions At*1 - At*5 (dossier 156-6, lot 2)}

\page{7}
Je voudrais maintenant faire passer les axiomes de disjonction-compatibilité \add{\uncertain{moins élégants}} avant ceux sur les lieux, pour ne pas les en faire dépendre. Les axiomes sont « facultatifs » pour un atelier, et il y aura pas mal d'implications mutuelles avec les axiomes sur les lieux.

Une notion essentielle ici est celle de « $F$ et $G$ \emph{disjointes} »
\[
  F \parallel G \overset{\mathrm{déf}}{\Longleftrightarrow} F \between G \text{ et } F \cap G = \varnothing_{\mathfrak{F}}
\]
notion qui ne dépend que de $\leq$. On a d'autre part
\[
  F \parallel G \Longleftrightarrow \forall\, X \in \widetilde{F},\ Y \in \widetilde{G},\ X \parallel Y
\]
d'où interprétation en termes de $X \parallel Y$ dans $\mathcal{M}$, qui y signifie
\[
  X \underset{\mathcal{M}}{\parallel} Y \Longleftrightarrow X \underset{\mathcal{M}}{\between} Y \text{ et } \lbrace X, Y \rbrace \text{ n'est pas minoré dans } \mathcal{M}, \leq \text{, i.e.\ } \nexists\, Z \text{ avec } Z \leq X,\ Z \leq Y
\]
(NB Il résultera des axiomes \add{précédents} que \ill{} \uncertain{indique} que $\nexists\, Z$ avec $Z \ll X$, $Z \ll Y$, \struck{i.e.\ \ill{}} et la réciproque \uncertain{résultera} d'axiome ci-dessous…)

\emph{At DC 1} (\emph{Axiome des ateliers spécieux}) Soient $X, Y \in \mathcal{M}$, \add{$K$} $= X \cap Y = \mathop{\mathrm{Sup}}_{Z \in \widetilde{X} \cap \widetilde{Y}} Z$. Supposons que pour \add{tous} $X', Y' \in \mathcal{M}$
\[
  X' \ll X,\ X' \parallel K,\ Y' \ll Y,\ Y' \parallel K \Longrightarrow X' \parallel Y' .
\]
Alors [NB il \ill{} dans tous les cas… cf.\ Cor] $X \between Y$.
\note{dans la ligne centrale, un signe est biffé devant $X'$ ; « cf.\ Cor » est écrit sous le crochet}

\emph{Cor} \quad Soient $F, G \in \mathfrak{F}$, $K = F \cap G = \mathop{\mathrm{Sup}}_{Z \in \widetilde{F} \cap \widetilde{G}} Z$. Alors $F \between G$ (ssi) $\forall\, X, Y \in \mathcal{M}$,
\[
  X \ll F,\ Y \ll G,\ X \parallel K,\ Y \parallel K \Longrightarrow X \parallel Y .
\]
\note{« ssi » (lecture incertaine) est cerclé ; le bas de la page, depuis « At DC 1 », est barré de trois longs traits obliques}

\marginal{Un atelier sera dit \emph{spécieux} s'il satisfait \uncertain{cette} condition \ill{} \uncertain{contre-ex.} : les lieux \ill{} \ill{}…}
\note{note oblique dans la marge gauche, en face de At DC 1, coupée par l'un des traits de biffure ; lue en partie}

\page{8}
J'explicite, à cause de son importance, la propriété suivante, bien qu'elle sera conséquence de diverses façons \ill{}

\struck{At DC 1 : Soient $X, Y \in \mathcal{M}$ avec $X \between Y$} \add{avec $X' \parallel X \cap Y = L$, i.e.\ $X' \parallel Z$ pour $Z \leq X, Y$}

\struck{At} Soient $X' \ll X$, \struck{$Y' \ll Y \Longrightarrow X' \parallel Y'$} $Y' \ll Y$ avec $Y' \parallel L$. Alors $X' \parallel Y'$

\emph{Cor 1} \quad Soient $F, G \in \mathfrak{F}$ avec $F \parallel G$. Alors $F' \ll F$, $G' \ll G \Longrightarrow F' \parallel G'$.

\emph{Cor.\ 2} \quad Soient $F, G \in \mathfrak{F}$, $F \between G$ \struck{\ill{}}, \add{$L = F \cap G$}, $F' \ll F$, $G' \ll G$. Supposons $F' \parallel L$, $G' \parallel L$, alors $F' \parallel G'$.
\note{depuis « At DC 1 » jusqu'ici, le haut de la page est barré de deux longs traits obliques ; le premier énoncé a d'abord été écrit avec $X \between Y$ puis récrit, l'hypothèse $X' \parallel X \cap Y = L$ étant ajoutée au-dessus de la ligne}

\marginal{NB Il \ill{} choisi \ill{} $X' \between Y'$, alors $X' \cap Y' = \varnothing$ \uncertain{immédiat} (\ill{} $Z \leq X'$, $Z \leq Y'$ $\Rightarrow$ $Z \ll X$, $Z \ll Y$, d'où $Z \ll X \cap Y = L$ ; $Z \leq X' \parallel L \Rightarrow Z \parallel L$ ; et $Z \ll L$ et $Z \parallel L$ \ill{} \ill{} $\Rightarrow$ $Z \mathrel{\overline{\parallel}} L = \varnothing$ !}
\note{note oblique dans la marge gauche, en face de ces énoncés, écrite ligne à ligne en formules ; l'ordre et plusieurs signes sont incertains}

\struck{\emph{At DC 2} (Axiome des ateliers spécieux) Soient $X, Y \in \mathcal{M}$, $K = X \cap Y = \mathop{\mathrm{Sup}}_{Z \in \widetilde{X} \cap \widetilde{Y}} Z$. Supposons que pour tous $X', Y'$, avec $X' \ll X$, $Y' \ll Y$, $X' \parallel K$, $Y' \parallel K$, on ait $X' \parallel Y'$. Alors}
\note{énoncé resté sans conclusion, encadré puis barré de traits obliques ; la ligne « Soient … Supposons que » est biffée à part}

On voudrait la

\emph{Proposition 1} \quad Soient $F, G \in \mathfrak{F}$, $F \between G$, $L = F \cap G$. Soient $F' \ll F$, $G' \ll G$. Alors
\[
  F' \between G' \Longleftrightarrow F'_{L} \between G'_{L}
\]
\struck{cf.\ GF VI, p.\ 73 — On a en tous cas $\Longrightarrow$. En effet,} \add{\ill{} \ill{}} Ceci est équivalent \ill{} (\ill{} \ill{}) à l'ancien At 12, qui prend le nom
\note{la Proposition reprend celle de GF VI, p.\ 73-74, où la même équivalence ($F' \between G'$ ssi $F'|K \between G'|K$) était ramenée à un « Cor.\ Main » puis posée en axiome At 12 (dossier 156-6, lot 4)}

\emph{At CD 1} \quad Soient $X, Y$ \add{ou \uncertain{ancien}} et $X', Y' \in \mathcal{M}$, $X \between Y$, $X' \ll X$, $Y' \ll Y$. Posons $K = X \wedge Y = \mathop{\mathrm{Sup}}_{Z \in \widetilde{X} \cap \widetilde{Y}} Z$. Alors $X'_{L} \between Y'_{L}$ (NB l'inverse \ill{} toujours vrai) implique (cf.\ GF VI) $X' \between Y'$
\note{« At CD 1 » est écrit sur un premier « At DD 1 » (ou « At DC 1 ») ; l'énoncé pose $K$ et conclut sur l'indice $L$, tel quel. Un trait vertical le borde à gauche}

\emph{Cor.\ Main 1} (cf.\ la « propriété » ci-dessus) \quad $F \parallel G$, $F' \ll F$, $G' \ll G \Longrightarrow F' \parallel G'$

\emph{Cor.\ Main 2} \quad Si $F \between G$, $F' \ll F$, $G' \ll G$, alors
\[
  F' \parallel G' \Longleftrightarrow F'_{L} \parallel G'_{L}
\]
\note{devant l'équivalence, deux mots sont lus « pour que » sans certitude}

\marginal{Critère de compatibilité \struck{\ill{}} \add{\ill{}}. Mais c'est faux dans le cas simplicial, ex : [dessin]. Mais on aura $X'_{L} \parallel Y'_{L} \Longrightarrow X' \mathrel{\overline{\parallel}} Y'$ (c'est un contre-ex.\ à At 11, et \ill{} At 11 p.\ 67, 68)}
\note{longue note oblique dans la marge gauche, en face de At CD 1 et des corollaires. Le dessin est un losange contenant un cercle coupé par un diamètre vertical, avec de petites croix autour ; il n'est pas redessiné. Le renvoi « p.\ 67, 68 » est à GF VI, où At 11 est le principe de recollement des raffinements (dossier 156-6, lot 4). Le signe barré de la conclusion est lu $\overline{\parallel}$ sans certitude}

\page{9}
Forme affaiblie de la Propriété 1

\struck{At CD 1' Soient $X, Y, X$}

\emph{Propriété $1'_{0}$} \quad $F, G, F', G'$ avec $F \between G$, $F' \ll F$, $G' \ll G$. Posons $L = F \cap G$, \struck{et supposons $F'_{L} \parallel G'_{L}$} Alors
\[
  F' \parallel G' \Longleftrightarrow F'_{L} \parallel G'_{L}
\]
\note{l'indice de « $1'_{0}$ » (ici et dans At CD $1'_{0}$) est un petit « o » sous le prime ; lecture incertaine. Un trait vertical borde l'énoncé à gauche}

L'implication $\Longrightarrow$ \uncertain{toujours valable}. \add{La validité de} l'implication \uncertain{en sens inverse} équivaut à :

\emph{At CD $1'_{0}$} \quad Soient $X, Y, X', Y' \in \mathcal{M}$, $X \between Y$, $X' \ll X$, $Y' \ll Y$. Posons $L = X \cap Y$ et supposons que $X'_{L} \parallel Y'_{L}$ (i.e.\ $\forall\, X'', Y'' \in \mathcal{M}$, $X'' \leq X'$, $Y'' \leq Y'$, $X'' \ll L$, $Y'' \ll L$, on ait $X'' \parallel Y''$) Alors $X' \parallel Y'$ (NB \uncertain{réciproque} \ill{} $=$ \uncertain{sans} At CD 1')

\marginal{Critère de disjonction \ill{} des raffinements de sous-figures}
\note{note oblique dans la marge gauche, soulignée, en face de At CD $1'_{0}$, que borde un double trait vertical}

\emph{Cor.\ 1} \quad Validité de Propriété 1'

\emph{Cor.\ 2} \quad Si $F, G, F' \in \mathfrak{F}$, avec $F \between G$, $F' \ll F$, supposant $F'_{L} = \varnothing_{\mathfrak{F}}$ (où $L = F \cap G$), alors $F'$ est disjoint de tout raffinement $G'$ de $G$.
\note{dans « $F, G, F'$ », le $F'$ est écrit sur un autre signe}

\emph{Cor.\ 3} \quad Soient $F \parallel G$, alors $F' \ll F$, $G' \ll G \Longrightarrow F' \parallel G'$.

\marginal{Corollaire 4 \quad Soient $F' \ll F$, $L \leq F$, $F'_{L} = \varnothing_{\mathfrak{F}}$ \ill{} $\Rightarrow$ $F' \parallel L$}
\marginal{Pour que $F'_{L} = \varnothing_{\mathfrak{F}}$, il est n.\ et s.\ que $F' \parallel L$. Cela il suffit vraiment \ill{} \ill{}, i.e.\ $F'_{L} = \varnothing_{\mathfrak{F}} \Rightarrow F' \parallel L$ (\ill{} « \ill{} » \ill{} l'axiome \ill{})}
\note{deux notes obliques dans la marge gauche, en face des corollaires ; lecture en partie incertaine}

J'ai tendance à regarder At CD comme « axiome obligatoire » pour les ateliers.

\struck{At CD 2 (Axiome des ateliers spécieux)}

\emph{Propriété 2} (des « ateliers spécieux »). Soient $F, G \in \mathfrak{F}$, \struck{\ill{}} on ne suppose pas $F \between G$, mais soit
\[
  L = F \cap G = \mathop{\mathrm{Sup}}_{Z \in \widetilde{F} \cap \widetilde{G}} Z .
\]
\struck{A} Supposons que pour tout $F' \ll F$ tel que $F'_{L} = \varnothing_{\mathfrak{F}}$ (i.e.\ $F' \parallel L$ \add{par At CD 1'}) et $G' \ll G$ tel que $G'_{L} = \varnothing_{\mathfrak{F}}$ (i.e.\ $G' \parallel L$) on ait $F' \parallel G'$. Alors $F \between G$.
\note{un trait vertical borde l'énoncé à gauche ; « par At CD 1' » est ajouté au-dessus de « (i.e.\ $F' \parallel L$) »}

\page{10}
Pour ceci, il \struck{faudrait} \add{faut} qu'on ait l'axiome suivant

\emph{At CD 2} \quad Soient $X, Y \in \mathcal{M}$, \struck{on} on ne suppose pas $X \between Y$ compatibles mais on introduit $K = X \cap Y = \mathop{\mathrm{Sup}}_{Z \in \widetilde{X} \cap \widetilde{Y}} Z$. On suppose que pour tout $X', Y' \in \mathcal{M}$ avec $X' \ll X$, $X'_{L} = \varnothing_{\mathfrak{F}}$ et $Y' \ll Y$, $Y'_{L} = \varnothing_{\mathfrak{F}}$, on ait $X' \parallel Y'$ \struck{\ill{}}. Alors $X \between Y$.
\note{l'énoncé introduit $K$ et emploie l'indice $L$, tel quel ; le $\parallel$ de la conclusion est repassé à l'encre}

\marginal{Critère de compatibilité des ateliers spécieux}

Cet axiome n'est pas satisfait \add{le plus souvent} pour les ateliers dits « \emph{polyédraux} », i.e.\ satisfaisant : \struck{pour tout} \add{deux} multistrates, \add{compatibles, minorées} \struck{\ill{}} \struck{disjointes} ont pour intersection une multistrate \add{i.e.\ $X \between Y$,} si $X, Y$ \struck{\ill{}} dans $\mathcal{M}$ \add{(par \uncertain{ex.}\ majorés)} \struck{et} minorés dans $\mathcal{M}$, alors \struck{\ill{} dans $\mathcal{M}$} $X, Y$ ont un inf dans $\mathcal{M}$, i.e.\ $X \wedge_{\mathfrak{F}} Y \in \mathcal{M}$. Pour ceux-ci,
\note{les ajouts et ratures de cette phrase sont interlinéaires ; l'ordre de lecture n'est pas sûr}

\emph{At CD 2'} \quad Soient $X, Y \in \mathcal{M}$, $K = X \cap Y = \mathop{\mathrm{Sup}}_{Z \in \widetilde{X} \cap \widetilde{Y}} Z$. On suppose que $\forall\, X', Y' \in \mathcal{M}$, $X' \ll X$, $X'_{L} = \varnothing_{\mathfrak{F}}$, $Y' \ll Y$, $Y'_{L} = \varnothing_{\mathfrak{F}}$, on ait $X' \parallel Y'$. \struck{Alors $X \between Y$.} Supposons \add{de plus} $K \in \mathcal{M}$, i.e.\ c'est un inf \struck{\ill{}} de $X, Y$ dans $\mathcal{M}$. Alors $X \between Y$.

\marginal{Critère de compatibilité pour ateliers polyédraux}

\emph{NB} \struck{Chacune des deux conditions implique At CD $1_{0}$.}

\emph{N.B.} \quad Tous les ateliers que j'ai en vue sont intuitivement satisfaisant At CD $1_{0}$, et l'un des deux axiomes At CD 2 ou At CD 2'.
\note{la ligne précédente est barrée de traits obliques ; celle-ci est bordée à gauche d'un double trait vertical, « l'un des deux axiomes » étant ajouté au-dessus de la ligne}

L'intérêt des \add{axiomes} At CD, c'est de rendre \uncertain{superflue} \add{\ill{}} $\between$ dans $\mathcal{M}$, contrôlée par les figures, i.e.\ \struck{la propriété \ill{} de $\leq$} dans $\mathcal{M}$, \struck{\ill{}} \uncertain{tel est} le rôle crucial de cette notion de disjonction \add{et} dans la construction d'ateliers à partir de ($\mathcal{M}, \leq, \ll$ \struck{$\between$}), \struck{\ill{}} ou maintenant en termes de $(\mathcal{M}, \leq, \ll, \parallel_{\mathcal{M}})$.

\marginal{Attention \quad La relation $\parallel$ dans $\mathcal{M}$ ne résulte pas \uncertain{comme} dans $\mathfrak{F}$ \ill{} de la seule donnée de $\leq$ — ici il faut \ill{} voir (en l'absence de $\between$) \ill{} \ill{} \uncertain{primitive} \ill{}}
\note{note oblique dans la marge gauche, au bas de la page, lue en partie ; le « $\mathcal{M}$ » de la première ligne est souligné}

\page{11}
Il serait donc intéressant de dégager les axiomes pertinents des ateliers spécieux ou polyédraux, en termes cette fois de $(\mathcal{M}, \leq, \ll, \parallel)$ — en laissant de côté la donnée de $\widetilde{\mathfrak{F}}$ (qui \uncertain{est} explicite dans le cas des \add{ateliers à} figures de type fini..), et dans le cas général la donnée de $\widetilde{\mathfrak{F}}$ se réduit à celle de certains antifiltres de $\mathfrak{F}$ — i.e.\ d'une notion de familles \emph{loc.\ finies} de strates..)

\struck{Dans le cas des ateliers spécieux, on a}

\struck{Axiome(s) des \ill{}}

\struck{\ill{}} \struck{On voudrait que les}
\note{ces deux lignes sont barrées de cinq traits obliques}

J'\uncertain{hésite} sur un \uncertain{bon} \add{\uncertain{moyen}} \struck{\ill{}}, dans la formalisation $\ll$ dans le cas d'un atelier simplicial \add{(plus exactement polyédral)}. J'ai admis que $\ll$ entre simplexes (d'un espace topologique, ou modéré, disons) était la relation $\ll$ des figures \emph{ensemblistes}, de sorte que j'ai accepté l'inclusion d'un $\Delta^{1}$ dans un $\Delta^{2}$
\note{suivent deux petits dessins : un triangle dont la base porte deux points joints par un arc intérieur, et un triangle aplati dont un côté est doublé d'un arc ; entre eux, « ou »}

ce qui me conduit à accepter des subdivisions de $\Delta^{2}$ en deux triangles comme ceci
\note{dessin : un triangle aux trois sommets marqués d'un point, coupé par un arc joignant les deux sommets de la base, avec une petite croix au milieu de la base}

ce que \uncertain{personne} ne regarderait comme une subdivision simpliciale ! C'est la notion des \emph{simplexes affines} dans un espace affine qui a guidé les \uncertain{yeux}, et il faut utiliser

\page{12}
ici la notion de raffinement \emph{\ill{}}. Dans GF V p.\ 19, 20, il fallait supposer, \struck{\ill{}} pour $F \overset{s}{\ll} G$, dans a) \struck{que} [au sens de $\forall\, A \in F$, $G^{A}$ a un plus petit élément] que pour $A \in F$, $B \in G$, $A \cap B$ soit ou bien vide, ou une \struck{face de} sous-strate de $A$ : En termes de la \ill{} $G$ sur $|F|$, i.e.\ une famille \uncertain{induite} par $G$ sur $|F|$, les $B \in G$ \ill{} vérifient que les traces sur $|F|$ des \ill{} $\overset{\circ}{B} \cap |F| \neq \varnothing$). \struck{Cela signifie ceci \ill{}}
\note{le renvoi est au chapitre V (dossier 156-5), p.\ 19-20 ; le petit « s » est écrit au-dessus du signe $\ll$. La fin du passage, depuis « En termes de », est barrée de cinq traits obliques et suivie d'une ligne biffée}

Il faudra faire une \uncertain{opération} \ill{} \uncertain{séparée} pour le critère de compatibilité (\ill{} « compatibilité stricte », ou « polyédrale ») — qui signifie que \ill{} $F \cup G$ \add{figures ensemblistes}, \struck{familles} où $F$ et $G$ \ill{} \ill{} figures, mais que c'est une figure ensembliste \emph{polyédrale}.
\note{« figures ensemblistes » est écrit au-dessus de « familles », souligné, ce dernier mot étant biffé}

En résumé, \struck{pour} pour la formulation des axiomes spécieux, il faut distinguer \uncertain{entre} ceux relatifs aux « ateliers polyédraux », et ceux relatifs aux « ateliers spécieux ». Quand on en viendra aux plongements dans $\mathrm{Fig}(\mathcal{L})$, dans le premier cas il faudra poser (pour les ateliers polyédraux ensemblistes, i.e.\ pas \ill{} \uncertain{fondamentaux} aux ateliers ensemblistes polyédraux, qui n'existent \uncertain{pas} !) que
\[
  \mathrm{Multomb}(F) \overset{\mathrm{pol}}{\ll} \mathrm{Multomb}\, G \Longrightarrow F \ll G .
\]
\note{devant le premier « Multomb », une majuscule est biffée}

\marginal{\ill{} \ill{} ateliers GF VI \ill{} p.\ 52}
\note{note oblique au bas de la marge gauche ; GF VI, p.\ 52, ouvre la discussion des ateliers « ensemblistes » (At ens 1 ; dossier 156-6, lot 3)}

\page{13}
\emph{Définition} \quad Un ens.\ ordonné $I$ est dit \emph{polyédral} si deux éléments minorés ont un inf. \struck{(NB Si $I$ est fini, cela revient à ceci : \ill{} \ill{} $I_{\leq x}$ et \ill{})} \add{NB cela signifie que si} $x, y \in I$ sont tels que $I_{\leq x} \cap I_{\leq y} \neq \varnothing$, alors $I_{\leq x} \cap I_{\leq y}$ a un plus grand élément. Quand \struck{\ill{}} \add{\ill{}} est \uncertain{inductif} (\uncertain{par ex.}\ \uncertain{fini}), il revient au même de dire que \struck{\ill{}} $I_{\leq x} \cap I_{\leq y}$ est filtrant, \struck{\ill{}} ce qui signifie
\[
  \forall\, x, y, x', y' \in I \text{ tels que } x', y' \leq x, y ,\quad \exists\, z \in I \text{ tel que } x', y' \leq z \leq x, y .
\]
\note{la parenthèse biffée de la deuxième ligne est remplacée par l'ajout interlinéaire « NB cela signifie que si » ; dans la formule, un $x'$ est écrit sur un autre signe}

Sous cette forme, on voit que la condition est autoduale, et signifie donc aussi que deux éléments majorés de $I$ ont un \uncertain{plus grand élément}.

\marginal{Cela signifie aussi \ill{} \ill{} de deux parties fermées \emph{strictes} \ill{} \ill{} — \ill{} \ill{} !}
\note{longue note oblique dans la marge gauche, en face de la définition, lue seulement par fragments}

Donc \emph{\ill{}} Supposons que $\forall\, x, x' \in I$ tels que $x' \leq x$, $I_{x' \leq \cdot \leq x}$ soit inductif \struck{Alors} (p.ex.\ fini). Alors si $I$ est polyédral, $I^{\mathrm{op}}$ l'est, et l'inverse est vrai si les $I_{x' \leq \cdot \leq x}$ sont inductifs.

Si $\mathcal{A} = (\mathfrak{F}, \leq, \ll)$ est un atelier, $F \in \mathfrak{F}$, \struck{la figure} $F$ est dit \emph{polyédral} si $\widetilde{F}$ est polyédral, i.e.\ si deux strates de $F$ non disjointes ont comme intersection une strate de $F$. L'atelier $\mathcal{A}$ est dit \emph{polyédral} si \struck{toute} \add{ses} figures de $\mathcal{A}$ est polyédrale, i.e.\ si on a ceci (en termes de $\mathcal{M}$) :

\emph{At pol 1} \quad $\forall\, X, Y \in \mathcal{M}$, si $X \underset{\mathcal{M}}{\between} Y$ et $X$ et $Y$ non disjoints (i.e.\ $\lbrace X, Y \rbrace$ minoré dans $\mathcal{M}$) alors \struck{\ill{}} $\mathrm{Inf}^{\mathcal{M}, \leq}(X, Y)$ existe.

\marginal{Définition des ateliers polyédraux, ou condition de \uncertain{bonne} intersection}

Un ens.\ ordonné $I$ est dit \emph{simplicial}, s'il est polyédral, et si pour tout $x \in I$, $S = I_{\leq x}$ est un \emph{simplexe} \emph{combinatoire}, i.e.\ ($\mathfrak{P}^{*}(J) =$ ens.\ des parties non vides de $J$) \ill{} isomorphe à un ensemble ordonné $\mathfrak{P}^{*}(J)$, où $J$ est un ens.\ \emph{fini}. Pour un ens.\ ordonné $S$, il revient au même de dire \add{polyédral}, que $S$ est fini, que $S$ est un simplexe combinatoire, ou que $S$ a un plus grand élément, et que chaque élément \uncertain{non minimal} a au moins deux prédécesseurs. [Alors on voit par induction sur card $S$ que tout élément de $S$ est
\note{la fin de cette phrase se lit mal : la place de l'ajout « polyédral » et la suite des conditions sont incertaines}

\page{14}
le sup \add{de l'un des} des éléments minimaux qu'il majore, et si $J$ est l'ens.\ de ces éléments minimaux, on trouve $S \simeq \mathfrak{P}^{*}(J)$…) Une \struck{figure} \add{figure $F$} est dite \add{(combinatoirement)} \emph{simpliciale}, si $\widetilde{F}$ est simplicial — i.e.\ si elle est polyédrale, et si ses strates sont des simplexes combinatoires. Un \emph{atelier} est dit \add{(combinatoirement)} \emph{simplicial} si toutes ses figures sont simpliciales, i.e.\ il satisfait, en plus de At pol 1 ci-dessus, la condition

\emph{At simp 1} \quad $\forall\, X \in \mathcal{M}$, \struck{l'ens.\ $\widetilde{X} = \mathcal{M}_{\leq X}$ est un simplexe combinatoire, \ill{} dans $(\mathcal{M}, \leq)$, alors $X$ a au moins deux éléments prédécesseurs} \add{maximaux}, i.e.\ $\partial X \overset{\mathrm{déf}}{=} \lbrace Y \in \mathcal{M} \mid Y < X \rbrace$ est vide, ou bien ($X$ minimal dans $(\mathcal{M}, \leq)$) ou admet au moins deux éléments maximaux.
\note{la première rédaction, biffée sur trois lignes, est remplacée par la seconde ; « maximaux » est encadré avec la fin biffée ; au-dessus de « $\overset{\mathrm{déf}}{=}$ », « $\partial X$ a au moins deux éléments » est ajouté sans que sa place soit sûre. Dans « At simp 1 », le deuxième mot est lu sans certitude}

J'en viens aux axiomes complémentaires

\emph{At C} (\struck{Axiome} \add{Critère} de \emph{compatibilité des raffinements de sous-figures}) Soient $X, Y, X', Y'$ \add{$\in \mathcal{M}$} avec $X' \ll X$, $Y' \ll Y$, $X \between Y$. Soit $L = X \cap Y$ ($= \widetilde{X} \cap \widetilde{Y} = \lbrace Z \in \mathcal{M} \mid Z \leq X, Y \rbrace$), soient $X'_{L} = X'.L = \mathrm{Inf}^{\ll}(X, L) = \lbrace X'' \in \widetilde{X}' \mid X'' \ll L \rbrace$, et $Y'_{L} = Y'.L = \dots$, les raffinements de $L$ induits par $X', Y'$ respectivement. \struck{Alors} [On sait déjà que $X' \between Y' \Rightarrow X'_{L} \between Y'_{L}$] Alors
\[
  X'_{L} \between Y'_{L} \Longrightarrow X' \between Y' .
\]
\note{« $\mathrm{Inf}^{\ll}(X, L)$ » ainsi, pour $\mathrm{Inf}^{\ll}(X', L)$}

\marginal{Le \uncertain{formule} \ill{} de suite en termes de \uncertain{prop.} ?}
\note{note oblique dans la marge gauche, en face de At C}

Cet axiome vise à la validité de la

\emph{Proposition} \quad Soient $F, G, F', G' \in \mathfrak{F}$, avec $F' \ll F$, $G' \ll G$, $F \between G$ (i.e.\ $F, G$ sous-figures d'une même figure $H$). Soit $L = F \cap G$, $F'_{L} = F'.L$, $G'_{L} = G'.L$. Alors
\[
  F' \between G' \Longleftrightarrow F'_{L} \between G'_{L} .
\]
\note{At C et cette Proposition reprennent At CD 1 et la Proposition 1 de la page 8, et, au-delà, la Proposition et At 12 de GF VI, p.\ 73-76}

\page{15}
Au point de vue ensembliste, pour les supports des Multomb correspondants, on trouve la figure suivante (où $H = F \vee G$)

\begin{tikzcd}[column sep=small, row sep=small]
 & & \mathrm{Omb}\, H \arrow[dl, no head] \arrow[dr, no head] & & \\
 & \mathrm{Omb}\, F \arrow[dl, no head] \arrow[dr, no head] & & \mathrm{Omb}\, G \arrow[dl, no head] \arrow[dr, no head] & \\
 \mathrm{Omb}\, F' \arrow[dr, no head] & & \mathrm{Omb}\, L \arrow[dl, no head] \arrow[dr, no head] & & \mathrm{Omb}\, G' \arrow[dl, no head] \\
 & \mathrm{Omb}\, F'_{L} & & \mathrm{Omb}\, G'_{L} &
\end{tikzcd}

\note{sur la page, les termes sont reliés par de simples traits, sans flèches}

où chacun des trois carrés est cartésien. Itou pour les déploiements

\begin{tikzcd}[column sep=small, row sep=small]
 & & \widetilde{H} & & \\
 & \widetilde{F} \arrow[ur, hook] & & \widetilde{G} \arrow[ul, hook'] & \\
 & & \widetilde{L} \arrow[ul, hook'] \arrow[ur, hook] & & \widetilde{G}' \arrow[ul] \\
 \widetilde{F}' \arrow[uur] & & & & \\
 & \widetilde{F}'_{L} \arrow[ul, hook'] \arrow[uur] & & \widetilde{G}'_{L} \arrow[uul] \arrow[uur, hook] &
\end{tikzcd}

\note{les flèches à crochet (inclusions) et les flèches simples sont données comme elles sont tracées ; la disposition est celle de la page, en losange}

La compatibilité de $F'$ et $G'$ signifie que $\forall\, X' \in \widetilde{F}'$, $Y' \in \widetilde{G}'$, on a $X' \between Y'$ — et la proposition dit qu'il suffit de le vérifier pour les strates $X', Y'$ telles que $X' \in \widetilde{F}'_{L}$, $Y' \in \widetilde{G}'_{L}$, i.e.\ telles que \struck{\ill{}} $X' \ll L$, \add{$Y' \ll L$} \struck{\ill{}} (ce qui équivaut d'ailleurs à : \struck{$X' \ll G$,} $X' \ll G$, $Y' \ll F$). \struck{\ill{}} \emph{d'où} \uncertain{découle} la conséquence

\emph{At D} (\struck{Axiome} \add{Critère} de \emph{disjonction des raffinements de sous-figures}) Soient $X, Y, X', Y' \in \mathcal{M}$, avec $X' \ll X$, $Y' \ll Y$, $X \between Y$. Soit $L = X \cap Y$. [On sait déjà que $X' \parallel Y' \Rightarrow X'_{L} \parallel Y'_{L}$] Alors
\[
  X'_{L} \parallel Y'_{L} \Longrightarrow X' \parallel Y'
\]
\note{un trait vertical borde l'énoncé à gauche}

\page{16}
Visée : la validité de la

\emph{Proposition} \quad Données comme dans prop.\ précédente. Alors
\[
  F' \parallel G' \Longleftrightarrow F'_{L} \parallel G'_{L}
\]

\emph{Cor.\ Main 1} \add{(At $\mathrm{D}_{0}$)} \struck{Soient} $F, G, F', G'$ avec $F \parallel G$, $F' \ll F$, $G' \ll G$, alors $F' \parallel G'$.

\marginal{cf.\ \ill{} (\ill{} Cor.\ 2')}
\note{note oblique dans la marge gauche ; une accolade y réunit la Proposition et le Cor.\ Main 1}

\emph{At pol 1} (\struck{Axiome} \add{Critère} de \emph{compatibilité pour ateliers polyédraux}) Soient $X, Y \in \mathcal{M}$. On ne suppose pas $X, Y$ compatibles, on introduit pourtant \add{$L = X \cap Y =$} $\mathop{\mathrm{Sup}}_{Z \in \widetilde{X} \cap \widetilde{Y}} Z$, \add{\uncertain{figure telle}} que $\widetilde{L} = \widetilde{X} \cap \widetilde{Y}$. Supposons que pour tout $X', Y' \in \mathcal{M}$, avec $X' \ll X$, $Y' \ll Y$, \struck{\ill{}} $X'_{L} = \varnothing$, $Y'_{L} = \varnothing$, on ait
\[
  X' \parallel Y' .
\]
\struck{Alors} Supposons de plus que \add{$L = \varnothing$ ou} $L \in \mathcal{M}$, i.e.\ $\lbrace X, Y \rbrace$ \add{\uncertain{s'il est minoré}} \struck{\ill{}} a un Inf dans $(\mathcal{M}, \leq)$. Alors $X \between Y$.
\note{le numéro de « At pol 1 » est écrit en surcharge ; il reprend le nom de la condition de bonne intersection de la page 13 pour un autre énoncé. Dans la ligne de $X', Y'$, un passage d'environ deux mots est lourdement biffé}

\marginal{On ne veut pas cet axiome, dans \ill{} \ill{}. \ill{} il suffit \ill{} compatibilité — \ill{}. L'axiome \ill{} \uncertain{avec} At $\mathrm{D}_{0}$ \ill{}}
\note{longue note oblique dans la marge gauche, en face de At pol 1, lue en partie ; ses dernières lignes sont barrées de traits obliques}

Visée : la validité de la proposition

\emph{Prop.} (Moyennant At pol 1) Soient $F, G \in \mathfrak{F}$. \struck{\ill{}} Soit $L = F.G$. On ne suppose pas $F \between G$. \struck{Supp} Alors $F \between G$ \uncertain{ssi} les conditions suivantes sont satisfaites :
\begin{itemize}
  \item[(i)] $\forall\, X \in \widetilde{F}$, $Y \in \widetilde{G}$, \struck{\ill{}} tels que $\lbrace X, Y \rbrace$ soit minoré dans $\mathcal{M}$, $X \cap Y \in \mathcal{M}$, i.e.\ \struck{Supp.} $\widetilde{X} \cap \widetilde{Y}$ est vide ou a un plus grand élément
  \item[(ii)] $\forall$ \struck{$X, Y$} $X, Y \in \mathcal{M}$, \struck{$X$} avec $X \ll F$, $Y \ll G$, \struck{\ill{}} $X_{L} = \varnothing$, $Y_{L} = \varnothing$, on ait $X \parallel Y$
\end{itemize}
\note{« $L = F.G$ » ainsi, avec le point de l'intersection fine ; « ssi » est lu sans certitude}

\page{17}
\begin{tikzcd}[column sep=small]
 \mathrm{Omb}\, F \arrow[dr, hook] & & \mathrm{Omb}\, G \arrow[dl, hook'] \\
 & \mathrm{Omb}(F) \cap \mathrm{Omb}(G) \arrow[d, hook] & \\
 & \mathrm{Omb}(L) &
\end{tikzcd}

\note{schéma de sa main en haut de la page, sans texte qui l'introduise ; les traits se terminent par un petit crochet, rendu par des flèches à crochet dans le sens où ils sont tracés. Il se rapporte à la Proposition de la page 16, où $L = F.G$}

\emph{At spéc} (Critère de \emph{compatibilité pour ateliers dits spécieux}) Comme At pol, mais en supprimant la condition sur $L$ (\uncertain{$L = \varnothing$} ou $L \in \mathcal{M}$).

\emph{Proposition} \uncertain{correspondante} pour ateliers spécieux : comme prop.\ ci-dessus, en supprimant (i).

\emph{NB}
\[
  \begin{array}{l}
    \mathrm{At\ D} + \mathrm{At\ spéc} \Longrightarrow \\
    \mathrm{At\ D} + \mathrm{At\ pol} \Longrightarrow
  \end{array}
  \quad \mathrm{At\ C} \Longrightarrow \mathrm{At\ D} \Longrightarrow \mathrm{At\ D}_{0} .
\]
\note{sur la page, les deux premières implications convergent vers « At C » ; le « NB » est dans la marge gauche, en face}

\marginal{il suffit \ill{} ici d'avoir \emph{At D'}, forme affaiblie de At D, où on affirme seulement (pour $X'_{L} = \varnothing$, $Y'_{L} = \varnothing$) $\Longrightarrow X' \parallel Y'$ (NB At D' $\Longrightarrow$ At $\mathrm{D}_{0}$)}
\note{note oblique dans la marge gauche, reliée par une flèche aux deux implications « $\Longrightarrow$ At C »}

\emph{At fil} (Critère des Sup filtrants) \add{loc.\ finis} Soient $F \in \mathfrak{F}$, $(F_{i})_{i \in I}$ une famille filtrante croissante \add{(pour $\leq$)} de raffinements $F_{i} \ll F$, « localement finie dans $F$ » i.e.\ $\forall\, X \in \widetilde{F}$, la famille des $F_{i} | X$ est stationnaire. Alors $\mathop{\mathrm{Sup}} F_{i}$ existe dans $\mathfrak{F}, \leq$.
\note{« loc.\ finis » est ajouté au-dessus de la parenthèse ; « (pour $\leq$) » au-dessus de « croissante », un mot biffé dessous}

\marginal{On voudrait supposer $\ll$ \uncertain{aussi} ($i \leq j \Rightarrow F_{i} \ll F_{j}$) ? (et $F_{i} \leq F_{j}$) (et prendre $\mathrm{Sup}$ pour $\ll$ ??)}
\note{note oblique dans la marge gauche, en face de At fil}

\emph{Cor.\ Main 1} \struck{\ill{}} (Moyennant At C) Recollement des raffinements, quand $F = \mathop{\mathrm{Sup}} F_{i}$ (deux versions, suivant qu'on recolle sur des sous-figures de $F$, ou sur ses strates)
\note{le recollement des raffinements était l'axiome At 11 de GF VI, p.\ 67-70 ; il devient ici un corollaire}

\page{18}
\emph{Corollaire 2} \quad Énoncé analogue à At fil, mais avec $(F_{i})$ filtrants croissants pour $\ll$ et $\leq$.

\emph{At supp} (Axiome des supports d'un raffinement) Soient \struck{$F$,} $X \in \mathcal{M}$ \struck{\ill{}}, $G \ll X$. Supposons $\exists\, Y \in \mathcal{M}$ avec $Y \parallel G$, $Y \mathrel{\overline{\parallel}} X$. Alors \struck{on peut} il existe une telle $Y$ \add{(\uncertain{choisie})} avec $Y \ll X$.
\note{« $G$ » est écrit sur un $F$ ; un trait vertical borde l'énoncé à gauche}

\marginal{Peut-être y a-t-il lieu de renforcer : Pour $Y$ donnée, $Y \parallel G$, $Y \mathrel{\overline{\parallel}} X$, $\exists\, Y' \in \mathcal{M}$, avec $Y' \ll X$, telle que $Y' \parallel G$, $Y' \mathrel{\overline{\parallel}} Y$ \struck{(ou \ill{})}. Ceci est lié à la question de l'étude de l'\uncertain{ombre} « complémentaire » de $G$… \uncertain{Où} \uncertain{veut-on en} venir…}
\note{longue note oblique dans la marge gauche, en face de At supp et de la Prop.\ qui suit ; la conclusion « $Y' \ll X$ » est écrite sur un $F$}

\emph{Prop.} \quad Idem pour At supp, avec $X$ remplacé par une figure $F$ quelconque ($G \ll F$).

\emph{Question} : Quand $Y \parallel G$, $Y \mathrel{\overline{\parallel}} F$, peut-on trouver $Y' \ll Y$ avec $Y' \ll F$ ? (On trouverait \struck{alors} automatiquement $Y' \parallel G$, par At D.) (Cela \add{serait} \struck{\ill{}} \ill{} \uncertain{so}, le « \ill{} \ill{} ».) Mais \uncertain{moyennant} At L 1, \uncertain{notre} At 8) \ill{} \ill{} \ill{} \ill{} $Y$ est un lieu, \uncertain{disons} $x$, et \ill{} la \uncertain{première} question
\[
  x \mathrel{\overline{\parallel}} Y \Longleftrightarrow x \ll Y \quad \text{i.e.\ } x \in \mathrm{omb}(Y) \qquad (\Longleftarrow \text{ est évident})
\]
i.e.
\[
  x \parallel Y \Longleftrightarrow x \notin \mathrm{omb}(Y)
\]
\note{devant « $x \notin$ », un « $x \parallel$ » suivi d'un mot est biffé}
Mais ceci n'est guère vérifié que pour les ateliers ensemblistes. Notons pourtant

\emph{At ens} \struck{D} (\emph{Critère de disjonction ensembliste}) \struck{Pour que} Soient $X, Y \in \mathcal{M}$. Pour que $X \parallel Y$, il \struck{suffit} \add{(faut et)} suffit que \struck{\ill{}} $\mathrm{Omb}(X) \cap \mathrm{Omb}(Y) = \varnothing$, i.e.\ qu'il n'existe pas de $Z \in \mathcal{M}$ avec $Z \ll X$, $Z \ll Y$.
\note{« At ens » est souligné ; le mot qui le suit, biffé, semble « Disj » ou « D ». L'énoncé est encadré de parenthèses à gauche ; « (faut et) » est ajouté au-dessus de « suffit »}

\page{19}
\emph{Cor 1} \quad Itou pour figures $F, G$ :
\[
  F \parallel G \Longleftrightarrow \mathrm{Omb}\, F \cap \mathrm{Omb}\, G = \varnothing
\]

\emph{Cor 2} \quad Soient $x \in \mathcal{L}$, $F \in \mathfrak{F}$. Alors
\[
  x \parallel F \Longleftrightarrow x \not\ll F \quad \text{i.e.\ } x \notin \mathrm{omb}(F)
\]

\emph{Cor.\ 3} \quad Soient $x, y \in \mathcal{L}$, alors
\[
  x \parallel y \Longleftrightarrow x \neq y .
\]

\emph{NB} \quad On a trivialement
\[
  \mathrm{At\ ens\ D} \Longrightarrow \mathrm{At\ supp} .
\]
\note{au-dessus de la formule, un premier « At supp » est biffé}
On notera que At supp est vérifié dans les ateliers auxquels j'ai songé \uncertain{jusqu'à présent} pratiques, mais non At ens D.

\note{un trait horizontal, au milieu de la page, sépare ce qui suit}

J'en viens aux axiomes impliquant les lieux (At L 1, 2, 3)

\emph{At L 1} (\emph{Axiome de divisibilité}) Soit $X \in \mathcal{M}$, alors $\mathrm{omb}(X)^{\circ} \neq \varnothing$ i.e.\ $\exists\, x \in \mathcal{L}$ avec $x \mathrel{\overset{\circ}{\ll}} X$.
\note{un trait vertical borde l'énoncé à gauche ; dans la deuxième mouture, c'est l'axiome des lieux At 8, (1.52) $\mathrm{omb}(X)^{\circ} \neq \varnothing$ (GF VI, p.\ 45)}

Cet axiome d'\uncertain{apparence} \uncertain{anodine} implique pourtant \uncertain{(ce qui} signifie, \uncertain{\ill{}}) que les « espaces » ou « contrées » qu'on \uncertain{décrit} par des ateliers, sont « infiniment divisibles », et en particulier (quand ils ne sont pas discrets) sont infinis : c'est l'introduction du « \uncertain{l'infini} actuel » tout court !

\page{20}
\emph{At L 2} (Critère de \emph{disjonction par les lieux}) Soient $X, Y \in \mathcal{M}$. Alors
\[
  X \parallel Y \Longleftrightarrow \forall\, x \in \mathrm{omb}(X),\ y \in \mathrm{omb}\, Y,\ \text{on a } x \parallel y
\]
(À ne pas confondre avec $\mathrm{omb}(X) \cap \mathrm{omb}(Y) = \varnothing$ !)
\note{dans la deuxième mouture, ce critère était At 9 (GF VI, p.\ 61)}

\emph{Cor} \quad Soient $F, G \in \mathfrak{F}$. Alors
\[
  F \parallel G \Longleftrightarrow \forall\, x \in \mathrm{omb}(F),\ y \in \mathrm{omb}(G),\ \text{on a } x \parallel y
\]
\note{ce corollaire est écrit à gauche, en retrait sur la marge, entre At L 2 et At L 3}

\emph{At L 3} (\struck{Axiome} \add{Critère} de disjonction pour les lieux) Soient $X, Y \in \mathcal{M}$, \struck{\ill{}} $X \between Y$, $X \neq Y$. Alors
\[
  x \in \mathrm{omb}(X)^{\circ},\ y \in \mathrm{omb}(Y)^{\circ} \Longrightarrow x \parallel y
\]

\marginal{équivalent, moyennant L 2, à}
\note{note oblique dans la marge gauche, en face de At L 3}

Équivalent à

\emph{Cor} \quad Soit $F \in \mathfrak{F}$. Alors deux lieux appartenant \add{aux} intérieurs de deux strates \uncertain{distinctes} de $F$ sont disjoints.

\emph{NB} \quad Quand on admet At D, alors il suffit de poser At L 3 dans le cas $Y \leq X$.

Si on admet At L 2, alors \struck{At} At C (\struck{pol}) \add{spéc} prend la forme

\emph{At CL} (\uncertain{spéc}) (Critère de \emph{compatibilité par les lieux, pour ateliers} \struck{polyédraux} \add{spécieux}) Soient $X, Y \in \mathcal{M}$, $L = X \cap Y = \mathop{\mathrm{Sup}}_{Z \in \widetilde{X} \cap \widetilde{Y}} Z$. Alors pour que $X$ et $Y$ soient comp., \struck{il suffit} \add{il suffit que} $\forall\, x \in \mathrm{omb}(X) \smallsetminus \mathrm{omb}(L)$, $y \in \mathrm{omb}(Y) - \mathrm{omb}(L)$, on ait $x \parallel y$.
\note{l'énoncé est bordé à gauche d'un trait vertical ; « spécieux » est écrit au-dessus de « polyédraux » biffé, et le mot entre parenthèses après « At CL » semble « spéc » écrit sur « pol » ; l'ajout « il suffit que » est lu sans certitude}

\marginal{Vaut-il mieux poser : il \uncertain{faut} et il suffit ?}
\note{note oblique dans la marge gauche, en face de At CL}

En effet, \struck{la condi} dans l'énoncé de \struck{At spéc} \add{At spéc} \add{avec $X, Y, L$} : $X' \ll X$, $Y' \ll Y$, la condition $X'_{L} = \varnothing$, $Y'_{L} = \varnothing \Rightarrow X' \parallel Y'$, cette dernière condition s'explicite, par L 2, par
\note{« At spéc » est écrit dans la marge gauche, souligné, en face de « En effet » ; la page s'arrête sur « par », la suite au lot 2}
\end{document}
